\section{Reverse Nodes in Groups of K}
\label{sec:ll-reverse-k-group}

\problemstatement{Given a linked list and an integer $k$, reverse the nodes
in consecutive groups of $k$. If the final group has fewer than $k$ nodes,
leave it as is. Return the head.}

\begin{patternbox}
This is the iterative reversal applied \textbf{block by block}: for each
group of $k$ nodes, check that $k$ nodes remain, reverse that block in
place, and stitch it between the already-processed part and the rest. A
dummy node anchors the stitching.
\end{patternbox}

\subsection*{Count, reverse a block, reconnect}

Maintain a pointer to the node before the current group. First verify $k$
nodes remain (walk ahead $k$ steps); if not, stop. Reverse the $k$ nodes
with the standard save--flip--advance loop. Then reconnect: the previous
group's tail points to the new (reversed) block head, and the block's tail
points to the following node. Move the ``previous group tail'' pointer to
the block's tail and repeat.

\begin{ideabox}[Reverse in Place, Then Re-stitch the Seams]
The only complexity beyond plain reversal is bookkeeping at the two seams of
each block -- linking the prior tail to the new head and the new tail to the
remainder. A dummy node removes the special case for the very first block.
\end{ideabox}

\subsection*{Algorithm}

\begin{enumerate}
  \item Dummy before head; \texttt{groupPrev = dummy}.
  \item Loop: check $k$ nodes remain; if not, break. Reverse the next $k$
        nodes; reconnect the seams; advance \texttt{groupPrev} to the
        block's (new) tail.
  \item Return \texttt{dummy.next}.
\end{enumerate}

\subsection*{Dry run}

\dryrun{$1\to2\to3\to4\to5$, $k=2$.}

\begin{itemize}
  \item Reverse $[1,2]\to2\to1$; reverse $[3,4]\to4\to3$; $[5]$ left as is.
  \item Result $2\to1\to4\to3\to5$.
\end{itemize}

\subsection*{C++ implementation}

\begin{lstlisting}
#include <bits/stdc++.h>
using namespace std;

struct ListNode {
    int val; ListNode* next;
    ListNode(int v) : val(v), next(nullptr) {}
};

ListNode* build(const vector<int>& a) {
    ListNode dummy(0), *t = &dummy;
    for (int x : a) { t->next = new ListNode(x); t = t->next; }
    return dummy.next;
}

ListNode* reverseKGroup(ListNode* head, int k) {
    ListNode dummy(0); dummy.next = head;
    ListNode* groupPrev = &dummy;

    while (true) {
        ListNode* kth = groupPrev;
        for (int i = 0; i < k && kth; i++) kth = kth->next;   // k-th node
        if (!kth) break;                                       // fewer than k left

        ListNode* groupNext = kth->next;
        ListNode* prev = groupNext;
        ListNode* curr = groupPrev->next;
        while (curr != groupNext) {                            // reverse the block
            ListNode* nx = curr->next;
            curr->next = prev;
            prev = curr;
            curr = nx;
        }
        ListNode* newTail = groupPrev->next;                   // old head -> new tail
        groupPrev->next = kth;                                 // stitch to new head
        groupPrev = newTail;
    }
    return dummy.next;
}

void print(ListNode* h) { for (; h; h = h->next) cout << h->val << " "; cout << "\n"; }

int main() {
    print(reverseKGroup(build({1, 2, 3, 4, 5}), 2));   // 2 1 4 3 5
    return 0;
}
\end{lstlisting}

\begin{complexitybox}
\textbf{Time Complexity.} $O(n)$ -- each node is reversed once.
\textbf{Space Complexity.} $O(1)$.
\end{complexitybox}

\begin{takeawaybox}
Reverse-in-$k$-groups is block-wise iterative reversal with careful seam
stitching. The ``check $k$ remain, reverse, reconnect'' loop is the hard-
problem culmination of the plain reversal skill built earlier.
\end{takeawaybox}
